% !TEX root = sga4.5-en.tex

\chapter*{Introduction}

The purpose of this volume is to make it easier for the non-expert to use
$\ell$-adic cohomology. I hope that it will often allow the reader to avoid
recourse to the dense expos\'es of SGA 4 and SGA 5. It also contains some new
results.

The first expos\'e, written by J.F. Boutot, surveys SGA 4. It gives the main
results --- with minimal generality, often insufficient for applications ---
and an idea of their proof. For complete results, or detailed proofs, SGA 4
remains indispensable.

``\nameref{II}'' contains a completed proof of the trace formula for the
Frobenius endomorphism. The proof is the one given by Grothendieck in SGA 5,
pruned of every unnecessary detail. This Report should allow the user to
forget SGA 5, which may then be regarded as a series of digressions, some of
them very interesting. Its existence will make it possible to publish SGA 5
as it stands shortly. It is complemented by the expos\'e ``\nameref{VI}''
which explains how the trace formula makes it possible to study
trigonometric sums, and gives examples.

The other expos\'es are aimed at a more limited audience, and their style
reflects this. Expos\'e ``\nameref{III}'' is a ``modular'' generalisation of
the Report, based on the study in SGA 4 XVII 5.5 of symmetric powers.
Expos\'e ``\nameref{IV}'' defines this class in various contexts, and gives
the compatibility between intersections and cup products. Gathered in
``\nameref{V}'' are a few known results for which a reference was missing,
together with a few compatibilities. Expos\'e ``\nameref{VII}'' is new. In
particular, in cohomology without supports, it gives finiteness theorems
analogous to those known in cohomology with compact supports.

For more details on the expos\'es, I refer to their respective
introductions.

Finally, I thank J.L. Verdier for allowing me to reproduce here his notes
``\nameref{VIII}.'' I believe they remain very useful, and they had become
impossible to find.

In the internal references of this volume, expos\'es are cited by an
abbreviated title, indicated between [ ] in the table of contents.

\vspace{20pt}
Bures-sur-Yvette, 20 September 1976

\vspace{5pt}
Pierre Deligne
